\subsection*{Resultados rectores y secuencia de demostración}
\label{paper:resultados-rectores}
La exposición inicial anticipa la función de la constante de acoplamiento;
su construcción se efectúa después de definir las hojas de APP, la
orientación TRIT, los operadores TPK y las prolongaciones regionales.
La definición~\ref{lib01:definicion} especifica el registro ordenado, su
lectura posicional y la coordenatización en la que se representan. La
sección~\ref{subsec:k-terminal} construye la extracción terminal y distingue
la información visible de la información firmada conservada en el estado.

La función aritmética del registro se establece en la construcción de
\(\alpha\). La sección~\ref{subsec:alpha-publicaciones-completas} contiene
la representación analítica completa, su recurrencia en todos los grados,
la convergencia de sus truncamientos y la identificación de su raíz con la
precoordenada generada. Su dependencia de esa precoordenada forma parte de
la definición del lector. Esta construcción precede a la escala de acción,
a las coordenadas angulares y a las realizaciones de área y gravitación.

La función de reconstrucción se expresa mediante resultados distintos y
componibles. La proposición~\ref{lib01:marco-ciclico} demuestra que la
órbita cíclica del registro constituye una base; la
proposición~\ref{lib01:operadores} determina qué respuestas permiten
recuperar un operador. El teorema~\ref{lib04:teo-balance} establece la
conservación bilateral para dos isometrías. El
teorema~\ref{lib04:teo-archivo} construye el archivo de profundidad
ilimitada y su inversa continua. Finalmente, los
teoremas~\ref{lib05:teo-biyeccion} y~\ref{lib05:teo-lector} incorporan
el refinamiento aritmético y calculan la transformación de los observables.

Estos resultados hacen explícita la jerarquía del artículo: generación del
estado, extracción de la constante, realizaciones aritméticas e
incidenciales, conservación de formas y recuperación efectiva de la
información. Las fórmulas de reconstrucción especifican en cada caso el
registro que reciben; la composición conserva el significado de sus
dominios, sus escalas y sus orientaciones.

% BEGIN CENTRALIDAD_K_20260921
La selección conjunta de la constante queda desarrollada antes de su
reconstrucción algebraica: las ordenaciones del repertorio terminal producen
\(19\,446\) registros; la compatibilidad excepcional conserva \(79\), y la
compatibilidad temporal determina un único registro \(K\). Se conservan las
condiciones de ambos filtros y su dependencia de la historia enriquecida.
La sección~\ref{act21:centralidad-constante} articula esta determinación con
la clausura aritmética, las representaciones incidenciales y la recuperación
operatoria.
% END CENTRALIDAD_K_20260921
